
Neural network weights encode architectural priors and training dynamics. A ResNet's 25 million parameters capture batch normalization statistics, residual connection patterns, and learned convolution kernels that distinguish it from a Vision Transformer with patch embeddings and self-attention mechanisms. Weight-space learning treats model parameters as data, enabling applications such as model zoo search, neural architecture prediction, and anomaly detection without executing the models.

Processing weights as data requires handling variable layer shapes and high dimensionality while preserving architectural information. A standard approach flattens each layer into tokens for sequence models, but systematic validation of this tokenization strategy has been limited. Does layer-wise processing retain sufficient structure for property prediction, or are cross-layer dependencies critical? Architecture family classification serves as a validation task: if tokenization loses essential structure, even coarse-grained discrimination between ResNet and ViT should fail.

Prior work in weight-space learning has demonstrated feasibility through task-specific evaluations. Hyper-representations have been shown to transfer across meta-learning tasks, model stitching analyses indicate that layers can be examined semi-independently, and weight statistics have been correlated with model properties. However, these approaches lack controlled experiments comparing tokenization strategies---specifically, statistical aggregation versus learned token embeddings---independent of downstream task complexity.

This work addresses the question: How do different layer-wise processing approaches compare in preserving structural signal for architecture classification? The contribution is a systematic validation on the timm Model Zoo (100 pretrained vision models across ResNet, ViT, EfficientNet, and ConvNeXt families). The main findings are:

\begin{enumerate}
\item \textbf{Tokenization validation}: Layer-wise processing (flatten, pad, per-layer normalize) achieves 80\% accuracy on architecture family classification, significantly exceeding the 60\% threshold and 25\% random baseline.
\end{enumerate}

\begin{enumerate}
\item \textbf{Baseline comparison}: Simple per-layer statistical aggregation (mean, standard deviation, L2 norm) matches transformer-based token embedding performance (both 80\%) on small datasets, revealing that family-level patterns do not require cross-layer attention when architecture families exhibit strong layer-level separability.
\end{enumerate}

\begin{enumerate}
\item \textbf{Capacity-data mismatch}: The transformer (2M parameters) overfits on 70 training samples (100\% train, 73.33\% validation), while the baseline (200K parameters) generalizes better (98.57\% train, 86.67\% validation), yet both achieve identical test accuracy, indicating that model capacity calibration is critical in weight-space learning.
\end{enumerate}

These results validate layer-wise weight processing as a viable foundation while establishing that simple statistical baselines provide strong comparison points. Future work on larger datasets and harder tasks may reveal conditions under which cross-layer reasoning provides measurable benefits.
